\newpage
\section{Parallelization}
\label{sec:parallelization}

\subsection{Sequential section}
Ho deciso inizialmente di eseguire i test sul mio PC personale, piuttosto che sulla workstation dell'università o su Google Colab, per evitare limitazioni dovute a sessioni temporanee e risorse limitate. In questo modo ho avuto maggiore controllo sul processo di compilazione e sull'esecuzione dei test.\\

\noindent
Questa parte del report fa riferimento solo all'analisi eseguita tramite il mio PC personale.

\noindent
I parametri della simulazione utilizzati sono:

\begin{itemize}
    \item $nstep = 200$
    \item $ni = 15000$
    \item $nj = 15000$
    \item $size = ni \times nj = 900\,000\,000$
\end{itemize}

\noindent
Ho compilato il codice sequenziale presente in \texttt{src/cuda.c} using the following options:

\begin{verbatim}
    gcc -O3 src/cuda.c -o build/cuda
\end{verbatim}

\noindent
Il tempo di esecuzione migliore ottenuto ($37 s$) costituirà il riferimento sequenziale per il calcolo dello speedup nella versione GPU.

\subsection{CUDA}

\subsubsection{Code}
La parallelizzazione è stata realizzata utilizzando CUDA per eseguire il calcolo dei passi temporali su GPU. Il kernel principale è \texttt{step\_kernel\_mod}, che replica il comportamento della funzione sequenziale \texttt{step\_kernel\_ref}.\\

\noindent
Il codice si può trovare in \texttt{src/cuda.cu}

\paragraph{Struttura del kernel:}
\begin{itemize}
    \item Ogni thread gestisce un punto della griglia 2D.
    \item Gli indici globali $i$ e $j$ vengono calcolati a partire dai \texttt{threadIdx} e \texttt{blockIdx}.
\end{itemize}

\paragraph{Gestione della memoria:}
\begin{itemize}
    \item Allocazione della memoria host e device.
    \item Copia dei dati iniziali dal host alla device.
    \item Copia dei risultati dalla device al host al termine della simulazione.
\end{itemize}

\paragraph{Controllo errori:}
\begin{itemize}
    \item Dopo ogni chiamata al kernel viene eseguito \texttt{cudaGetLastError()}.
    \item Viene calcolato l'errore massimo rispetto al riferimento CPU per verificare la correttezza.
\end{itemize}

\begin{lstlisting}
__global__ void step_kernel_mod(int ni, int nj, float fact, float* temp_in, float* temp_out)
{
  int i00, im10, ip10, i0m1, i0p1;
  float d2tdx2, d2tdy2;

  int i = threadIdx.x + blockIdx.x * blockDim.x;
  int j = threadIdx.y + blockIdx.y * blockDim.y;

  if(i > 0 && i < ni-1 && j > 0 && j < nj-1){
    // find indices into linear memory
    // for central point and neighbours
    i00 = I2D(ni, i, j);
    im10 = I2D(ni, i-1, j);
    ip10 = I2D(ni, i+1, j);
    i0m1 = I2D(ni, i, j-1);
    i0p1 = I2D(ni, i, j+1);

    // evaluate derivatives
    d2tdx2 = temp_in[im10]-2*temp_in[i00]+temp_in[ip10];
    d2tdy2 = temp_in[i0m1]-2*temp_in[i00]+temp_in[i0p1];

    // update temperatures
    temp_out[i00] = temp_in[i00]+fact*(d2tdx2 + d2tdy2);
  }
}
\end{lstlisting}



\subsubsection{Compilation}
La compilazione CUDA è stata effettuata con il comando:

\begin{verbatim}
    nvcc -O3 -arch=sm_50 cuda.cu -o cuda.exe -DDEBUG=0
\end{verbatim}

\textbf{Significato dei flag:}
\begin{itemize}
    \item \texttt{-O3}: ottimizzazione massima.
    \item \texttt{-arch=sm\_50}: architettura target della GPU.
    \item \texttt{-DDEBUG=0}: disabilita il debug per migliorare le performance.
\end{itemize}

\subsubsection{Execution}
Per il kernel CUDA, la scelta di \texttt{threadsPerBlock} e \texttt{numBlocks} è stata la seguente:

\begin{lstlisting}
    dim3 threadsPerBlock(16, 16);
    dim3 numBlocks((ni + 15) / 16, (nj + 15) / 16);

    // ---- CUDA events start ----
    cudaEvent_t start, stop;
    cudaEventCreate(&start);
    cudaEventCreate(&stop);

    // Execute the modified version using same data
    printf("Execute GPU kernel...\n");
    cudaEventRecord(start, 0); // start timing

    for (istep = 0; istep < nstep; istep++) {
        step_kernel_mod <<<numBlocks, threadsPerBlock>>> (ni, nj, tfac, d_temp1, d_temp2);
        cudaDeviceSynchronize();

        cudaError_t err = cudaGetLastError();
        if (err != cudaSuccess)
            printf("CUDA Error: %s\n", cudaGetErrorString(err));
        
        float* d_tmp = d_temp1;
        d_temp1 = d_temp2;
        d_temp2 = d_tmp;
    }

    cudaEventRecord(stop, 0); // stop timing
    cudaEventSynchronize(stop);

\end{lstlisting}

\noindent
Ogni blocco di 16x16 thread calcola i punti della griglia, mentre i blocchi sono organizzati per coprire l'intera matrice.  

\noindent
La temporizzazione è stata eseguita tramite \texttt{cudaEvent\_t} per misurare il tempo del kernel GPU.

\subsubsection{Results}

Profiling (command: \texttt{nvprof build/cuda})

\paragraph{950M:}
{\scriptsize
\begin{verbatim}
Execute the CPU-only reference version...
- Elapsed time: 0.188 s
Execute GPU kernel...
- Elapsed time: 140.515 ms
Max Error = 0.000008
The Max Error of 0.00001 is within ACCEPTABLE bounds.
- Total elapsed time: 2.460 s
==768== Profiling application: build/cuda
==768== Profiling result:
Type  Time(%)      Time     Calls       Avg       Min       Max  Name
GPU:   72.95%  202.16ms         2  101.08ms  2.7665ms  199.40ms  [CUDA memcpy HtoD]
       24.93%  69.101ms       200  345.50us  342.59us  355.94us  step_kernel_mod(int, int, float, float*, float*)
        2.12%  5.8774ms         2  2.9387ms  2.5962ms  3.2812ms  [CUDA memcpy DtoH]
API:   90.99%  1.69655s         2  848.27ms  699.40us  1.69585s  cudaMalloc
        7.10%  132.39ms       200  661.97us  423.90us  7.6760ms  cudaDeviceSynchronize
        0.96%  17.932ms         4  4.4831ms  3.1292ms  7.7131ms  cudaMemcpy
        0.45%  8.4117ms         2  4.2059ms  2.4000us  8.4093ms  cudaEventCreate
        0.43%  8.0093ms       200  40.046us  9.9000us  1.0583ms  cudaLaunchKernel
        0.04%  796.50us         2  398.25us  302.30us  494.20us  cudaFree
        0.01%  213.70us       200  1.0680us     200ns  31.300us  cudaGetLastError
        0.00%  90.700us         1  90.700us  90.700us  90.700us  cudaEventSynchronize
        0.00%  81.700us       101     808ns     200ns  31.300us  cuDeviceGetAttribute
        0.00%  30.100us         2  15.050us  10.700us  19.400us  cudaEventRecord
        0.00%  5.7000us         2  2.8500us  1.3000us  4.4000us  cudaEventDestroy
        0.00%  4.3000us         1  4.3000us  4.3000us  4.3000us  cuDeviceGetPCIBusId
        0.00%  3.9000us         1  3.9000us  3.9000us  3.9000us  cudaEventElapsedTime
        0.00%  2.3000us         3     766ns     400ns  1.4000us  cuDeviceGetCount
        0.00%  1.6000us         2     800ns     300ns  1.3000us  cuDeviceGet
        0.00%  1.5000us         1  1.5000us  1.5000us  1.5000us  cuDeviceGetName
        0.00%     800ns         1     800ns     800ns     800ns  cuDeviceTotalMem
        0.00%     400ns         1     400ns     400ns     400ns  cuDeviceGetUuid
\end{verbatim}
}

\paragraph{T4:}
{\scriptsize
\begin{verbatim}
Execute the CPU-only reference version...
- Elapsed time: 36.761 s
Execute GPU kernel...
- Elapsed time: 1741.090 ms
Max Error = 0.000015
The Max Error of 0.00002 is within ACCEPTABLE bounds.
- Total elapsed time: 46.074 s
==5152== Profiling application: /content/drive/MyDrive/Colab Notebooks/HPC-CUDA/build/cuda
==5152== Profiling result:
Type  Time(%)      Time     Calls       Avg       Min       Max  Name
GPU:   67.30%  1.73103s       200  8.6551ms  8.0181ms  11.542ms  step_kernel_mod(int, int, float, float*, float*)
       16.79%  431.89ms         2  215.95ms  203.45ms  228.45ms  [CUDA memcpy HtoD]
       15.91%  409.13ms         2  204.56ms  201.12ms  208.01ms  [CUDA memcpy DtoH]
API:   62.34%  1.73679s       200  8.6840ms  8.0240ms  11.548ms  cudaDeviceSynchronize
       30.23%  842.24ms         4  210.56ms  201.46ms  228.69ms  cudaMemcpy
        7.16%  199.38ms         2  99.688ms  144.08us  199.23ms  cudaMalloc
        0.13%  3.7350ms       200  18.674us  5.4010us  429.51us  cudaLaunchKernel
        0.13%  3.6890ms         2  1.8445ms  870.86us  2.8181ms  cudaFree
        0.01%  168.47us       114  1.4770us     185ns  65.476us  cuDeviceGetAttribute
        0.00%  67.106us       200     335ns     105ns     674ns  cudaGetLastError
        0.00%  33.409us         2  16.704us  10.078us  23.331us  cudaEventRecord
        0.00%  27.898us         2  13.949us  1.1280us  26.770us  cudaEventCreate
        0.00%  12.296us         1  12.296us  12.296us  12.296us  cuDeviceGetName
        0.00%  6.8260us         1  6.8260us  6.8260us  6.8260us  cuDeviceGetPCIBusId
        0.00%  4.4120us         1  4.4120us  4.4120us  4.4120us  cudaEventSynchronize
        0.00%  3.6190us         1  3.6190us  3.6190us  3.6190us  cudaEventElapsedTime
        0.00%  3.3080us         2  1.6540us  1.1110us  2.1970us  cudaEventDestroy
        0.00%  1.6790us         3     559ns     233ns  1.2100us  cuDeviceGetCount
        0.00%     716ns         2     358ns     179ns     537ns  cuDeviceGet
        0.00%     529ns         1     529ns     529ns     529ns  cuDeviceTotalMem
        0.00%     372ns         1     372ns     372ns     372ns  cuModuleGetLoadingMode
        0.00%     302ns         1     302ns     302ns     302ns  cuDeviceGetUuid
\end{verbatim}
}

\noindent
I risultati della simulazione GPU sono stati confrontati con la versione CPU e lo speedup è calcolato come: 
\[
\text{speedup} = \frac{T_\text{CPU}}{T_\text{GPU}}
\]
\begin{table}[h!]
\centering
\begin{tabular}{lccc}
\hline
\textbf{GPU} & $\mathbf{T_{CPU}}$ (ms) & $\mathbf{T_{GPU}}$ (ms) & \textbf{Speedup} \\
\hline
950M & 34967 & 15086 & x2.32 \\
T4   & 37178 & 1645  & x22.60 \\
\hline
\end{tabular}
\caption{Confronto tra tempi CPU e GPU e speedup}
\end{table}